%% =====================================================================
%%  Chapter 1: Introduction
%%  Background -> problem -> objective -> overview of the approach
%%  -> contributions -> organization of the thesis.
%%  A reader should understand why the work matters and what you did
%%  from this chapter alone.
%% =====================================================================
\chapter{Introduction}
\label{chap:introduction}

\section{Background}
\label{sec:background}

\begin{writingnote}
Start from the broad social or academic context and narrow it down to the problem studied in this thesis.
Support every claim with evidence and cite the literature where appropriate.
Briefly explain the key terms, such as those in the title, when you first use them;
give their precise definitions in the Preliminaries chapter.

%% Start a new paragraph with an empty line. Do not use \\ to break lines.
In the next paragraph, briefly mention what the main existing studies have achieved,
and then explain the problem that they have not yet solved.
Leave the detailed review of existing work to the Related Work chapter.
Concrete examples and numbers help the reader see why the problem needs to be solved.

%% The following books are useful for learning how to structure and write a thesis
%% (they also serve as an example of citing two works at once).
For advice on how to structure and write a thesis, see~\cite{zobel2014writing,glasmandeal2010science}.
\end{writingnote}

In recent years, XX has posed a challenge in \ldots.
Existing approaches \ldots, which motivates an investigation of \ldots.

\section{Objective}
\label{sec:objective}

\begin{writingnote}
State the objective of this thesis clearly in one or two sentences.
Make it correspond to the problem described in the previous section,
and make it concrete enough that the Conclusion chapter can judge whether it was achieved.
Then briefly describe the approach and what the experiments demonstrate.
Do not assume superiority: state the performance, applicability and limitations supported by the evidence.
Revise this section after you have written the other chapters, and check that the objective and the contributions
match what the Proposed Method and Computational Experiments chapters show.
\end{writingnote}

The objective of this thesis is to investigate \ldots.
The contributions of this thesis are as follows:
\begin{itemize}
  \item We formulate the XX problem as a YY.
  \item We propose a method for the YY problem using XX.
  \item We establish the performance, applicability and limitations of the proposed method through computational experiments.
\end{itemize}

\section{Organization of the Thesis}
\label{sec:organization}

%% Give chapters, sections, figures, tables and equations a name with \label{...}
%% and refer to them with \cref{...}. Words such as "Chapter", "Figure" and
%% "Eq." and the numbers are inserted automatically.
The rest of this thesis is organized as follows.
\Cref{chap:preliminaries} defines the problem studied in this thesis and describes the existing method on which the proposed method is based.
\Cref{chap:related-work} reviews related work and positions this study.
\Cref{chap:method} describes the proposed method.
\Cref{chap:experiments} presents and discusses the results of computational experiments.
\Cref{chap:conclusion} concludes the thesis and outlines future work.

%% Writing tips
%%  - Use the present tense for general facts and for describing your method,
%%    and the past tense for what you did in the experiments.
%%  - "We" is acceptable in theses in this field.
%%  - Define every abbreviation the first time you use it, e.g.,
%%    "mixed-integer programming (MIP)".
%%  - Avoid contractions (write "do not" instead of "don't").
